% ==============================================================================
% 4. SCHEDULE, TEAM ORGANIZATION, AND RISK CONTROL
% ==============================================================================
\newpage
\section{Schedule, Team Organization, and Risk Control}

\subsection{Planned schedule}

\begin{table}[H]
\centering
\caption{Project milestones and required evidence}
\label{tab:schedule}
\begin{tabularx}{\textwidth}{>{\raggedright\arraybackslash}p{0.18\linewidth}>{\raggedright\arraybackslash}p{0.36\linewidth}L}
\toprule
\textbf{Milestone} & \textbf{Primary work} & \textbf{Required evidence} \\
\midrule
Week 3 &
Finalize the problem, scope, and annotation guide; conduct the Steps 1--2 pilot; initialize the web/API scaffold with sample records. &
Proposal; pilot and feasibility logs; interface skeleton; unresolved-item list. \\
By Week 6 &
Complete collection, annotation, preprocessing, EDA, and data splitting. &
Image manifest; annotation guide; EDA; frozen split lists. \\
Weeks 7--10 &
Train and tune at least three model families; select the deployment model on validation; freeze all candidates; evaluate all models and the majority baseline once on the same test; integrate the selected model into the API. &
Logs and frozen configurations; shared-test metric table; confusion matrices; selected model; offline/API parity check. \\
Reporting week &
Complete cloud deployment, demonstration, report, and slides. &
PDF report; source-code archive; saved model; slides; public URL. \\
\bottomrule
\end{tabularx}
\end{table}

\subsection{Team organization}

Each member will lead one workstream:
\begin{enumerate}
  \item Problem definition and source verification
  \item Annotation, EDA, and data splitting
  \item Baseline and model development
  \item Web and API development
  \item Integration, deployment, and reporting
\end{enumerate}

All members will cross-check outputs, and at least two will annotate independently. The model lead cannot change test labels after viewing predictions. Final assignments will be agreed by the team.

\newpage
\subsection{Risk register}

\begin{table}[H]
\centering
\caption{Principal risks and planned responses}
\label{tab:risks}
\begin{tabularx}{\textwidth}{>{\raggedright\arraybackslash}p{0.30\linewidth}J}
\toprule
\textbf{Risk} & \textbf{Planned response} \\
\midrule
Too few valid images or severe class imbalance &
Test feasibility in the pilot; add streets, self-captured images, and minority-class cases where needed. Report the actual achieved scale and never substitute near duplicates for independent observations. \\
Ambiguous labels or data leakage &
Revise the guide, use independent adjudication, and group by street, keeping team capture sessions, parent images, and near-duplicate clusters together. Use only verified-street images in the main experiment and keep the test set sealed. \\
External API access or usage restrictions &
Verify access, licence conditions, and attribution during the pilot. Prefer team-captured images when an external source is unsuitable; review personal information and keep a compliant offline demonstration set. \\
Model overfitting or weak street transfer &
Monitor validation learning curves, keep street-disjoint splits, limit tuning to validation data, and report subgroup and error analyses. \\
Deployment or external-service outage &
Pin dependencies, record the model version, validate the offline/API parity, and prepare a fallback demonstration that does not fabricate live data. \\
\bottomrule
\end{tabularx}
\end{table}

\subsection{Expected deliverables and conclusion}

The final submission is expected to contain a verified image manifest, annotation guide, EDA and experiment notebooks, reproducible configurations, trained-model artefacts, a report and presentation, source code for the web interface and API, and a working deployment URL. Claims will remain proportional to the evidence: the map will describe recorded observations, evaluation will describe the sampled streets and conditions, and the classifier will be presented as a survey-support tool rather than a legal or accessibility certification system. Observations will be reported at street level, without identifying individual vendors or people.
